\begin{tikzpicture}[x=0.6cm,y=0.6cm, every node/.style={font=\small}]
  % Eingabe 5x5: Zeilen "0 0 9 9 9"
  \foreach \r in {0,...,4}{
    \foreach \c/\v/\g in {0/0/0,1/0/0,2/9/27,3/9/27,4/9/27}{
      \node[draw=black!50, minimum size=0.6cm, inner sep=0pt, fill=black!\g] at (\c,-\r) {\v};
    }
  }
  % Hervorgehobener Ausschnitt (Zeilen/Spalten 1-3 -> Index 0..2)
  \draw[figblue, line width=1.6pt] (-0.5,0.5) rectangle (2.5,-2.5);
  \node[figlabel] at (2,1.2) {Bildausschnitt (5$\times$5)};

  \node at (5.1,-1) {$\ast$};
  % Kernel
  \foreach \r in {0,1,2}{
    \foreach \c/\v in {0/1,1/0,2/-1}{
      \node[draw=figorange, fill=figorange!15, minimum size=0.6cm, inner sep=0pt] at (6.2+\c,-0.0-\r-0) {\v};
    }
  }
  \node[figlabel] at (7.2,1.2) {Filter (3$\times$3)};
  \node at (10.3,-1) {$=$};
  % Ausgabe 3x3: je Zeile -27 -27 0
  \foreach \r in {0,1,2}{
    \foreach \c/\v in {0/-27,1/-27,2/0}{
      \node[draw=black!50, minimum size=0.6cm, inner sep=0pt, fill=figgreen!15] at (11.4+\c,-\r) {\v};
    }
  }
  \draw[figblue, line width=1.6pt] (10.9,0.5) rectangle (11.9,-0.5);
  \node[figlabel] at (12.4,1.2) {Ergebnis (3$\times$3)};
  \node[figlabel, text width=7cm] at (6.5,-5.6) {Beispiel: $3\cdot(0\cdot1+0\cdot0+9\cdot(-1)) = -27$\\(großer Betrag bei einer senkrechten Kante)};
\end{tikzpicture}
